% =========================================================
% Proof: Uniqueness of Infimum
% Source: volume-iii/analysis/bounding/notes/notes-bounds-extremal-values.tex
% =========================================================

\subsection*{Uniqueness of Infimum}
\label{prf:inf-is-unique}

\begin{remark*}[Return]
$\leftarrow$ \hyperref[prop:sup-unique]{Back to Proposition in Notes}
\end{remark*}

\begin{proof}
\Claim Let $A \subseteq \mathbb{R}$ be nonempty and bounded below.
If $s$ and $t$ are both infima of $A$, then $s = t$.

\Given $s = \inf A$ and $t = \inf A$.
\Goal $s = t$.

Since $s = \inf A$, $s$ is a lower bound of $A$.
Since $t = \inf A$ is the greatest lower bound, $s \le t$.

Since $t = \inf A$, $t$ is a lower bound of $A$.
Since $s = \inf A$ is the greatest lower bound, $t \le s$.

\Therefore $s \le t$ and $t \le s$.
By antisymmetry of the order on $\mathbb{R}$, $s = t$. \AsReq
\end{proof}

\begin{remark*}[Proof shape]
Mirror image of the supremum uniqueness proof: apply the greatest-lower-bound
condition in each direction to derive $s \le t$ and $t \le s$, then conclude
by antisymmetry.
\end{remark*}

\begin{remark*}[Dependencies]
Definition of infimum (greatest lower bound); antisymmetry of $\le$ on $\mathbb{R}$.
\end{remark*}
